\documentclass[exercice,a4]{thedocument}%
\usepackage[french]{babel}%
\usepackage{textcomp}%
\usepackage{pstricks}%
\usepackage[np]{numprint}%
\startdatakeys
    \source{21GENMATNC1}
    \BO{2026}
    \cycle{4}
    \competencesDuSocle{CA,RA}
    \objectifsApprentissage{C4A3aa,C4D2ae,C4A1aa}
    \motsCles{pourcentage,fraction,équation,boule,volume,trigonométrie,angle}
\enddatakeys
\begin{document}%
\def\degres{\textsuperscript{\textdegree}}%

\medskip

Pour chaque affirmation répondre par vrai ou faux. Justifier chaque réponse.

\medskip

\textbf{Affirmation 1 :} 50\,\% de \np{10350} c'est \np{10300}.

\medskip

\textbf{Affirmation 2 :} $\dfrac73$ est la forme irréductible de $\dfrac{42}{18}$.

\medskip

\textbf{Affirmation 3 :} L'équation $2x - 4 = -x + 5$ a pour solution $3$.

\medskip

\textbf{Affirmation 4 :} L'arrondi à l'unité près du volume d'une boule de diamètre $21,6$~cm est \np{42213}~cm$^3$.

\smallskip

\emph{On rappelle la formule du volume d'une boule } $V = \dfrac43 \pi R^3$.

\medskip

\begin{minipage}{0.52\linewidth}
\textbf{Affirmation 5 :} Dans la figure codée ci-contre, la  mesure de l'angle $\widehat{\text{DNB}}$, arrondie à l'unité près, est $18\degres$.
\end{minipage}\hfill
\begin{minipage}{0.38\linewidth}
\psset{unit=1cm}
\begin{pspicture}(5.5,2.4)
\pspolygon(0.4,0.2)(5.1,0.2)(0.4,1.9)%BND
\psframe(0.4,0.2)(0.6,0.4)
\uput[dl](0.4,0.2){B} \uput[dr](5.1,0.2){N} \uput[u](0.4,1.9){D} \uput[l](0.4,1.05){4 cm} \uput[d](2.75,0.2){12 cm}
\end{pspicture}
\end{minipage}

\bigskip

\textbf{Affirmation 6 :} On peut composer 6 codes différents avec un cadenas à 3 chiffres qui respecte les conditions suivantes :

\begin{itemize}
\item[$\bullet~~$]les deux premiers chiffres sont choisis parmi 1 ; 2 et 3 ;
\item[$\bullet~~$]un chiffre peut apparaître deux fois;
\item[$\bullet~~$]le dernier chiffre est 6.
\end{itemize}

\bigskip

\showthesource
\end{document}
